% =====================================================
\chapter{Conclusions and Future Directions}
\label{chap:conclusion}

% =====================================================
\section{Conclusions}
\label{sec:conclusions}

\noindent This thesis studied the problem of atmospheric image synthesis for the purpose of generating realistic environmental domain shifts for computer vision systems. The motivation behind this research stems from a fundamental challenge of modern visual recognition models: although remarkable performance has been achieved under controlled conditions, systems are often degraded significantly when applied in adverse environmental settings. Fog, haze, sandstorms, snowfall, and underwater degradation are some atmospheric phenomena that change visibility, contrast, colour distributions, and texture statistics of scenes, resulting in complex distribution shifts that are still difficult to handle with conventional training paradigms. Collecting atmospheric datasets at a large scale is also expensive and logistically difficult, particularly for rare or hazardous environmental conditions. These observations motivate the development of methods for synthetic atmospheric data generation that can generate realistic and diverse weather-induced domain shifts.

\noindent The main goal of this thesis was to investigate whether recent advances in diffusion-based generative models can be leveraged to generate realistic atmospheric datasets that are not only visually plausible, but also induce meaningful representation-level changes in modern vision systems. Instead of concentrating on a single approach to atmospheric generation, the research investigated multiple complementary perspectives of atmospheric synthesis, progressing from purely appearance-based methods to physically motivated image formation approaches.

\noindent In the first stage of the research, atmospheric generation was investigated using diffusion-based style transfer. Modern image generation models have shown remarkable capabilities for artistic style transfer, illumination condition modification, and scene appearance change while preserving semantic content. The thesis built on these advances to explore whether atmospheric conditions could be treated as transferable visual styles. Multiple approaches were evaluated, including InstantStyle, StyleShot, and other diffusion-based frameworks, across several atmospheric settings. The proposed method for this stage --- \textbf{V0 (InstantStyle with CSD Guidance)} --- combined block-selective IP-Adapter injection with a Contrastive Style Descriptor latent correction loop, achieving the lowest CLIP confidence (0.514 ViT-L-14) of any style transfer method evaluated. Experiments showed that diffusion models are very capable of replicating weather-specific visual features such as colour shifts, contrast attenuation, and environmental textures. The atmospheric changes produced by these methods were often visually convincing and showed a great degree of stylistic agreement with the weather reference images.

\noindent The style-transfer experiments, however, also revealed fundamental limitations. While atmospheric appearance could be successfully reproduced, the induced effects were often physically inconsistent. Visibility loss was generally uniform across the entire image, independent of the distance of objects from the camera. Foreground objects remained unrealistically sharp, and the atmospheric effect resembled an overlay rather than a genuine environmental effect. These observations suggested that atmospheric conditions cannot be fully characterised by image style alone. Unlike artistic styles, atmospheric effects are sensitive to scene geometry, depth, light transport, and the physics of the environment. Purely appearance-based transfer methods, although effective at visual modification, turned out to be insufficient for modelling the underlying mechanisms of atmospheric image formation.

\noindent Inspired by these findings, the second phase of the research investigated optimization-guided diffusion as a means to explicitly include atmospheric constraints during the generation process. Atmospheric synthesis was reformulated as a constrained optimization problem, with style-aware objectives embedded directly into the reverse diffusion path. The work involved Contrastive Style Descriptors, Constrained Stable Diffusion, Langevin Dynamics, and Proximal Langevin Dynamics. These experiments yielded useful insights into the controllability of diffusion models and showed that atmospheric appearance can be guided by explicit optimization objectives rather than relying solely on feed-forward conditioning.

\noindent The optimization-guided experiments uncovered significant findings. First, they showed that atmospheric style can be imposed throughout the generation process, not only at initialization. Second, they emphasised the importance of guidance scheduling: optimizing in later denoising steps generally leads to a better trade-off between style fidelity and content preservation. Third, important limitations of embedding-space similarity objectives were identified. Semantic style embeddings do not fully capture the unique frequency and perceptual characteristics of atmospheric phenomena. Investigation of perceptual and frequency-aware objectives confirmed a more general conclusion: realistic atmospheric synthesis requires representations that incorporate physically meaningful image properties, not only visual similarity.

\noindent Optimization-guided diffusion improved controllability and enabled deeper insight into the atmospheric generation process, but also revealed the increasing complexity of trying to learn atmospheric physics implicitly through generative models. This realization motivated the final stage of the thesis. Instead of training diffusion models to learn the entire atmospheric image formation process from data, the research investigated the integration of explicit physical models into the synthesis pipeline.

\noindent This transition resulted in the development of the physics-grounded atmospheric synthesis framework presented in Chapter~\ref{chap:phase2_pipeline}. The proposed architecture was derived from knowledge acquired in prior experimental studies. The framework does not treat atmospheric generation as a single monolithic task, but breaks the problem into three independent stages: semantic scene generation, atmospheric appearance transfer, and physical image formation. Clean agricultural scenes are first generated with SDXL~\cite{podell2023sdxlimprovinglatentdiffusion}. Atmospheric characteristics are then introduced with InstantStyle~\cite{wang2024instantstyle} and IP-Adapter~\cite{ye2023ipadapterfidelityadapters} guidance. Finally, physically meaningful environmental effects are enforced with depth-aware atmospheric compositing based on the Koschmieder scattering model~\cite{koschmieder1924theorie}.

\noindent One of the most important contributions of this thesis is the recognition that semantic content generation and atmospheric generation should be treated as separate problems. Previous architectures were prone to semantic leakage, content drift, and instability when weather information and scene content were controlled jointly. The proposed framework decouples these processes, achieving substantially improved controllability by enabling atmospheric transformations to be applied without compromising scene structure or object identity. This architectural insight was key to the success of the final system.

\noindent Explicit atmospheric physics further altered the nature of the generated outputs. Unlike style-transfer approaches that only modify appearance, the physics-based compositor introduced depth-dependent visibility degradation, accumulation of atmospheric light, and weather-specific scattering behaviour. Refinements for sandstorms, snowfall, and snow accumulation demonstrated that different atmospheric phenomena have distinct physical characteristics that cannot be adequately represented by generic image transformations. The inclusion of weather-specific physical models was therefore an important factor in achieving realistic atmospheric synthesis.

\noindent The quantitative results obtained throughout the thesis offer strong evidence for the effectiveness of this approach. CLIP-based confidence analysis showed consistent decreases in classifier confidence when evaluated on generated atmospheric images, confirming significant distribution shifts at the representation level. \textbf{V9 (CLIP content subtraction + extended block injection + Koschmieder compositing) achieved the lowest confidence of any method across the entire study (0.228 ViT-L-14 / 0.227 ViT-H-14)}, establishing it as the best overall configuration and the most effective method for creating out-of-distribution atmospheric training data. Linear probing experiments demonstrated that the generated datasets contain meaningful atmospheric structure that can support downstream learning tasks. The distribution-specificity tests revealed an important observation: atmospheric adaptation is highly dependent on the properties of the generated distribution. This suggests that dataset generation methods can be as important for robustness under environmental shift as model architecture.

\noindent The results taken together support a broader conclusion about atmospheric image synthesis. Atmospheric generation should not be viewed purely as a style-transfer problem, nor purely as an optimization problem. Appearance, geometry, and physics work together to form realistic atmospheric synthesis. Diffusion models provide a powerful mechanism to generate semantically rich and visually plausible imagery, while explicit image formation models provide the physical structure necessary for realistic environmental transformations. The best atmospheric synthesis systems will likely integrate the advantages of both approaches rather than relying on one alone.

% =====================================================
\section{Future Work}
\label{sec:future_work}

\noindent The investigations presented in this thesis open several promising directions for future research. While the proposed atmospheric synthesis framework successfully integrates diffusion-based generation, style-guided transfer, and physics-grounded image formation, it also highlights opportunities to extend the realism, scalability, and applicability of synthetic atmospheric datasets.

\paragraph{Learned Atmospheric Compositors (KDA-Net and Beyond)}
The most immediate and directly motivated extension of this work is the development of a \emph{learned atmospheric compositor} to replace the hand-crafted Koschmieder stage. During the research process, an architecture named KDA-Net (Koschmieder-guided Depth-Aware Network) was prototyped: a lightweight encoder-decoder with three submodules --- a texture encoder, a haze-appearance head, and an image combiner --- trained with a combination of VGG perceptual loss and CSD contrastive loss (without adversarial training). The key insight was that classical Koschmieder compositing produces a \emph{globally smooth} depth-correlated veil that, while physically grounded, lacks the local texture irregularities present in real atmospheric images: patches of denser fog near water surfaces, scattering halos around light sources, and subject-specific depth-of-field effects. A learned compositor, trained on the output-to-reference pairs generated by the existing pipeline, could capture these sub-pixel details and produce higher-fidelity atmospheric synthesis without abandoning the depth-consistency guarantee of the physics model. The primary obstacle encountered was data volume: with only 57 real reference images in the evaluation set, the learned component was at risk of overfitting to their specific atmospheric textures. Future work with larger reference collections --- or with self-supervised objectives derived from the Koschmieder transmission law itself as a physics-regulariser --- could train a practical learned compositor that generalises across atmospheric conditions.

\paragraph{Style-Conditioned Physics Parameter Estimation}
A related but simpler learning-based direction is to replace the fixed global atmospheric parameters ($\beta$, $A$) derived from a single dark-channel prior estimate with a \emph{style-conditioned estimator} that infers per-region scattering density from the reference image. The current pipeline applies a single global $\beta$ and a single $A$ to every pixel, which approximates the atmospheric model accurately only for scenes without strong depth discontinuities. A small convolutional estimator trained to regress per-patch ($\beta$, $A$) values from style reference crops --- conditioned on depth-map statistics from DPT-Large --- could produce spatially varying atmospheric parameters and significantly improve depth-coherence in scenes with mixed near-field and far-field objects. This requires no adversarial training and could be trained on synthetic paired data generated by the existing pipeline, making it an accessible near-term extension.

\paragraph{Physics-Informed Diffusion Models}
One longer-term avenue is to build atmospheric physics directly into the diffusion process itself. In the current framework, physical image formation is a post-processing stage applied after atmospheric style transfer. This modular design provides flexibility and interpretability, but it means the diffusion model receives no physical feedback during generation --- it produces a stylised image without awareness of whether the resulting depth-attenuation profile is physically consistent. Future systems could benefit from incorporating Koschmieder-derived guidance terms directly into the reverse diffusion path: at each denoising step, a differentiable atmospheric rendering loss could nudge the predicted clean image towards configurations with depth-coherent haze. Such physics-constrained diffusion models would unify the style-transfer and compositing stages, potentially eliminating the colour-cast and content-leakage artefacts that motivated the block-selective injection and CLIP content subtraction techniques developed in this thesis.

\paragraph{Video and Temporal Consistency}
Another promising avenue is to extend the framework from image generation to video generation. In the real world, weather is inherently dynamic: fog density varies over time, snowfall has temporal structure, and sandstorms exhibit complex particle trajectories. The current framework operates on still images and does not enforce any temporal consistency. Future work could combine video diffusion models with physics-based atmospheric simulation to create temporally coherent weather sequences for testing autonomous systems in dynamic environments.

\paragraph{Compound Weather Conditions}
In this thesis, atmospheric phenomena were modelled separately, with distinct models for fog, haze, sandstorms, snowfall, and underwater environments. In practice, environmental conditions often coincide — fog with rain, snowstorms with high winds, or dust storms at sunset. Extending the framework to support multi-weather synthesis would increase realism and enable the construction of more challenging benchmark datasets.

\paragraph{Richer Atmospheric Physics}
The weather-specific refinements for sandstorms and snowstorms suggest a broader opportunity to incorporate richer atmospheric physics. The present implementation relies on simplified scattering models and heuristic accumulation mechanisms. Future work could include physically based rendering techniques, volumetric atmospheric simulation, radiative transfer models, and particle-based environmental modelling. Such approaches could substantially improve realism, especially under extreme weather conditions where multiple scattering effects are significant.

\paragraph{Adversarial Domain Shift Maximisation}
The V9 pipeline was designed and evaluated to maximise domain shift --- as measured by CLIP zero-shot confidence reduction --- using fixed pipeline parameters. A natural extension is to replace this fixed design with an adversarial optimisation loop: a \emph{generator} (the atmospheric pipeline parameterised by $\beta$, $A$, and IP-Adapter injection weights) and a \emph{discriminator} (a frozen pretrained CLIP or ViT encoder evaluating softmax confidence) are placed in competition, with the generator trained to minimise the discriminator's confidence on generated images while the discriminator remains frozen. This reduces to maximising the entropy of the classifier output over generated images --- an objective that could be optimised by differentiating through the Koschmieder compositor (which is differentiable with respect to $\beta$ and $A$) while treating the diffusion model as a black box sampled at each iteration. Preliminary experiments in this direction during the project showed that naive entropy maximisation leads to degenerate solutions (fully occluded grey images that trivially confuse any classifier); future work would need to combine entropy maximisation with a content-preservation regulariser and an image-quality prior.

\paragraph{Distribution-Aware Synthesis}
The distribution-specificity experiments in Chapter~\ref{chap:phase2_pipeline} showed that atmospheric adaptation is highly dependent on the characteristics of the generated distribution. This finding implies that dataset generation itself should be formulated as an optimisation problem. Instead of generating atmospheric images solely for visual realism, future systems could explicitly optimise synthetic datasets to maximise robustness gains for downstream tasks. Specifically, the per-class calibration failures observed in Experiments~3 and~5 --- where certain visually ambiguous categories (boar, cow) remained poorly handled even after LoRA adaptation --- suggest that a data-generation policy could be learned to over-sample the atmospheric conditions that most confuse the target model, closing the gap between visual fidelity and practical adaptation performance.

\paragraph{Broader Evaluation and Applications}
The current thesis primarily evaluated representation-level distribution shifts with vision-language models. Future work could evaluate the created datasets on a wider range of downstream tasks including object detection, semantic segmentation, instance segmentation, visual navigation, autonomous driving, agricultural monitoring, and robotic perception. Such studies would provide additional insight into the practical utility of synthetic atmospheric data across diverse application domains.

\paragraph{Domain Transfer}
The framework developed in this thesis was targeted at agricultural and rural environments. However, the underlying methodology is largely domain-independent. Future work could apply the same synthesis pipeline to urban scenes, aerial and satellite imagery, maritime environments, and indoor settings. These extensions would allow for the generation of diverse atmospheric datasets across a wide range of real-world applications.

\paragraph{Multimodal Foundation Models}
Finally, large multimodal foundation models are arriving rapidly and create new opportunities for atmospheric robustness research. Vision-language models, multimodal agents, and embodied AI systems are increasingly deployed under diverse environmental conditions. Synthetic atmospheric datasets created with frameworks such as the one proposed in this thesis could play an important role in understanding, evaluating, and improving the robustness of such next-generation systems.

\bigskip
\noindent In summary, this thesis has demonstrated that realistic atmospheric image synthesis can be achieved through a combination of diffusion-based generative modelling, optimization-based control, and physically grounded image formation principles. The research arc progresses from atmospheric style transfer through constrained diffusion to physics-aware atmospheric synthesis. The resulting framework offers a practical way to generate atmospheric datasets that produce significant environmental distribution shifts while maintaining visual realism and semantic consistency. It is hoped that the methods, insights, and findings presented in this thesis will contribute to future advances in atmospheric synthesis, domain adaptation, robustness evaluation, and the broader development of weather-aware computer vision systems.
